\section*{全册地图：主线与母题}
\addcontentsline{toc}{section}{全册地图：主线与母题}

先写课程的 1--3 条主线：本课程反复把哪些对象变成哪些工具？期末题最常考的母题是什么？

\begin{itemize}
  \item \textbf{主线一：} 用一句话说明知识如何串起来。
  \item \textbf{主线二：} 用一句话说明计算、证明或建模的核心动作。
  \item \textbf{综合关：} 说明哪些题会跨板块，需要哪些检查习惯。
\end{itemize}

\section{第一板块：示例知识板块}

\subsection{示例考点：方法名称}

\textbf{核心结论：} 写最小可执行理论、公式、条件。

\textbf{Use When：} 写题目出现哪些信号时使用本方法。

\textbf{标准套路：}
\begin{enumerate}
  \item 第一步：识别对象与条件。
  \item 第二步：选择定理、公式、变量替换或构造。
  \item 第三步：计算、化简，并做边界/符号/单位检查。
\end{enumerate}

\textbf{常见陷阱：} 写最容易丢分的点。

\textbf{真题接口：} \source{2023-2024, Final, 1}
